\section{股东资格与股东权利}
\begin{enumerate}
  \item 股东资格。
        \begin{enumerate}
          \item 成为股东对行为能力无要求。
          \item 股东资格确认之诉。公司为被告，其他利害关系人为第三人。
          \item 出资证明书。
          \item 股东名册。内部以名册为准，外部以登记为准。
        \end{enumerate}
  \item \zd{名义股东和实际股东}。每年必考。
        \begin{enumerate}
          \item 代持股。\zd{代持协议}。三重关系。
          \item 对内。际股东不能直接向公司主张权利。
          \item 外部。名义股东为准，对外补充责任。（是有权）处分股权 \yjt 类推善意取得。风险非常大。
          \item 浮出水面。过半数同意 \hz 半隐。
        \end{enumerate}
  \item 冒名股东。非法行为，没有合意。
  \item 权利。
        \begin{enumerate}
          \item 分红。
                \begin{enumerate}
                  \item 实缴比例 \hz \zd{全体}另有约定 100\% 同意。
                  \item 分红权之诉。载明具体分配方案的股东会决议。不交不支持。公司为被告。
                  \item 强制分红。滥用权力。
                  \item 时间。章程 \yjt 决议 \yjt 做出 1 年。
                \end{enumerate}
          \item （有限公司）优先认购。增资时，按照实缴出资比例优先认缴。
                \begin{enumerate}
                  \item \al{甲乙丙 100W，增资到 600，甲 0，乙 100，丙 200。甲 1/3 \yjt 1/6}。股权被稀释。
                  \item 实践中，间接把小股东踢出去；搭便车被踢出局。
                \end{enumerate}
          \item 知情权。非法目的除外。伏笔，后面可能跟着一连串诉讼。
                \begin{enumerate}
                  \item 书面，正当目的，无任何前置程序。
                  \item 一般文件能看能印，会计账簿只看不印。查账权之诉，能翻旧账。
                  \item 有限能印能看，股份只能看，没有查账权。
                  \item 不得实质剥夺。\qd{约定查账需要经过大股东同意}。不得放弃。
                  \item 可以带会计师、律师等负有保密义务的。要股东在场。
                \end{enumerate}
        \end{enumerate}
  \item 义务。不重要。
  \item 控股股东。超过 50\% \hz 足以对表决产生重大影响。
  \item 关联关系和关联交易。包括控股股东、实控人、董监高等。
  \item 股东的直接诉讼。股东受损，直接告董事、高管。侵权之诉。
  \item \zd{股东的代表诉讼/派生诉讼/代位诉讼}。公司受损，公司又不告。
        \begin{enumerate}
          \item 有限公司无要求，股份公司连续持股 180 天 + 1\%。
          \item 起诉时有股东资格，能溯及既往侵权。
          \item 用尽内部救济前提。书面交叉请求。
          \item 所得利益归公司，合理费用归个人。
          \item 调解需经董事会 \hz 股东会决议。
        \end{enumerate}
\end{enumerate}
